% !TeX root = ../../main.tex
\subsection{The negative-axis primitive and its coefficient bound}
\label{subsec:b:primitives}

We first fix the primitive geometrically at each genus. Uniformity
in the genus is a separate step: it uses a small interval near zero,
larger inverse disks on a compact negative interval, and growing
regular disks at negative infinity.

For integers $k\geq0$ and real $u>0$, the modified Bessel functions
in the following argument are defined by
\begin{equation}\label{eq:b:modified-bessel}
 I_k(2\sqrt u)=\sum_{a\geq0}\frac{u^{a+k/2}}{a!(a+k)!}.
\end{equation}
We use their elementary integral representation
\cite[(10.25.2), (10.32.3)]{DLMF}; the estimates needed from it
are included in the proof.

\begin{lmm}[Growing regular disks]\label{lmm:b:growing-disks}
For $u\geq1$ and integers $k\geq0$,
\begin{equation}\label{eq:b:modified-bounds}
 |(w')^{(k)}(-u)|=u^{-k/2}I_k(2\sqrt u)
                  \leq u^{-k/2}I_0(2\sqrt u),\qquad
 \tfrac13e^{\sqrt u}\leq I_0(2\sqrt u)\leq e^{2\sqrt u}.
\end{equation}
On $\mathbb D(-u,\sqrt u/10)$ the derivative $w'$ differs
from $I_0(2\sqrt u)$ by less than a quarter of that positive
number. Consequently, for each fixed $h\geq1$,
$\mathcal L_h(0;w(-u),y)$ is holomorphic for
\begin{equation}\label{eq:b:growing-radius}
 |y|<c_h\sqrt u\,I_0(2\sqrt u),
 \qquad |\mathcal L_h|\leq C_h I_0(2\sqrt u)^{h-1},
\end{equation}
with constants independent of $u$.
\end{lmm}
\begin{prf}
Termwise differentiation of $w'$ proves the equality in
\eqref{eq:b:modified-bounds}. For real $z>0$, the integral formula
is
\[
 I_k(z)=\frac1\pi\int_0^\pi e^{z\cos\varphi}
                                      \cos(k\varphi)\,\dd\varphi.
\]
Positivity follows from the series, and the integral gives
$I_k(z)\leq I_0(z)$. The upper exponential bound is immediate.
On $0\leq\varphi\leq\pi/3$, $\cos\varphi\geq1/2$, giving the
lower bound in \eqref{eq:b:modified-bounds}.
Taylor's formula now gives
\[
 |w'(-u+s)-I_0(2\sqrt u)|
       \leq I_0(2\sqrt u)(e^{|s|/\sqrt u}-1).
\]
For $|s|\leq\sqrt u/10$, the parenthesis is at most $1/9<1/4$;
one may use $e^t\leq(1-t)^{-1}$ for $0\leq t<1$.
Apply Lemma~\ref{lmm:b:regular-disk} with
$\beta=I_0(2\sqrt u)$ and $r=\sqrt u/10$ to conclude.
\end{prf}

\begin{lmm}[The geometric primitive]\label{lmm:b:primitive}
For integers $g\geq h\geq1$ with $g\geq2$,
\begin{equation}\label{eq:b:primitive}
 \mathcal G_{g,h}(0;0)
                     =\int_{-\infty}^0 L_{g,h}(0;x)\,\dd x.
\end{equation}
The integral is absolutely convergent at each such $g,h$. The
genus-one term $\mathcal G_{1,1}$ is excluded from this assertion.
\end{lmm}
\begin{prf}
On the negative axis write $x=w(-u)$, $u>0$. This parametrizes
the full axis: $w'(-u)=I_0(2\sqrt u)>0$, and its integral tends
to infinity by \eqref{eq:b:modified-bounds}.
In the finite formula \eqref{eq:s07:finite-jets},
$W_1=I_0(2\sqrt u)$ and
\[
 |W_d|\leq u^{-(d-1)/2}I_0(2\sqrt u)\qquad(d\geq2).
\]
Thus each term for genus $G\geq2$ is bounded by a constant times
$I_0(2\sqrt u)^{-2G+2}u^{-(2G-3+h)/2}$.
The derivative $\partial_x=-I_0(2\sqrt u)^{-1}\partial_u$
and \eqref{eq:b:modified-bounds} show by induction that, for every
fixed integer $m\geq0$,
\[
 |\partial_x^m\mathcal F_{G,h}(0;w(-u))|
             \leq C_{G,h,m}I_0(2\sqrt u)^{-2G+2-m}
                         \qquad(u\geq1).
\]
All powers of $u$ produced by differentiation are nonpositive and
may be discarded in this bound. The genus-one logarithm satisfies
the corresponding estimate for $m\geq1$:
$|\partial_x^m\log W_1|\leq C_m I_0(2\sqrt u)^{-m}$.

In the finite sine transformation for $\mathcal G_{g,h}$, a
genus-$G=g-k$ term has $2k$ derivatives. Its exponential factor
is therefore $I_0(2\sqrt u)^{-2g+2}$. A term with $G=1$
can occur only for $h=1$, and then $2k=2g-2\geq2$, so the
logarithm has already been differentiated. Hence
$\mathcal G_{g,h}(0;w(-u))\to0$ as $u\to\infty$.
An additional $x$ derivative gives an absolutely integrable bound
after multiplication by $|\dd x|=I_0(2\sqrt u)\,\dd u$.
The fundamental theorem of calculus proves \eqref{eq:b:primitive}.
For $g=h=1$, the undifferentiated expression is
$-\log I_0(2\sqrt u)/24$ and has no finite limit at infinity;
this explains the stated exception.
\end{prf}

The next elementary coefficient estimate is sufficient for the
primitive bound. The asymptotic transfer of a specified principal
term is a different statement and is used in Section 8.

\begin{lmm}[A two-endpoint coefficient bound]\label{lmm:b:coefficient-bound}
Let $\mathcal H(Y)$ be holomorphic in the unit disk and continue
across its circle except at $Y=\pm1$, with dented neighborhoods
there. Suppose on the circle and these neighborhoods that
\begin{equation}\label{eq:b:coefficient-hypothesis}
 \begin{aligned}
 |\mathcal H'(Y)|&\leq C
       (1+|1-Y|^{-1/2}+|1+Y|^{-1/2}),\\
 |\mathcal H''(Y)|&\leq C
       (1+|1-Y|^{-3/2}+|1+Y|^{-3/2}).
 \end{aligned}
\end{equation}
Then for integers $N\geq1$,
\begin{equation}\label{eq:b:coefficient-bound}
 |[Y^N]\mathcal H(Y)|\leq C' N^{-3/2},
\end{equation}
with $C'$ uniform for uniform hypotheses and angular margins.
\end{lmm}
\begin{prf}
The first derivative is integrable at each endpoint, so
$\mathcal H$ has matching continuous boundary values there.
Its circle restriction is absolutely continuous. Integrate its
Fourier coefficient once by parts, obtaining an outer factor
$1/N$ and the Fourier integral of the circle derivative
$\frac{\dd}{\dd\varphi}\mathcal H(e^{i\varphi})$.
Here $\varphi$ is a real circle parameter. On arcs of length
$1/N$ about the two endpoints, the integral of its absolute value is
$O(N^{-1/2})$, by the first bound in
\eqref{eq:b:coefficient-hypothesis}. On the complementary arcs,
integrate once more by parts. Their boundary terms are
$O(N^{-1/2})$, and the integral of its next derivative divided by $N$ is also
$O(N^{-1/2})$, since
$N^{-1}\int_{1/N}^1t^{-3/2}\,\dd t=O(N^{-1/2})$.
The original outer factor gives \eqref{eq:b:coefficient-bound}.
The Fourier representation follows either by radial limits in
Cauchy's formula or by the integrable derivative bounds just used.
Every estimate is uniform in the indicated constants.
\end{prf}

\begin{prp}[The unmarked primitive bound]\label{prp:b:primitive-bound}
For each fixed integer $h\geq1$ there is $C_h>0$ such that, for
$g\geq\max\{2,h\}$ and $N=2g+h-2$,
\begin{equation}\label{eq:b:primitive-bound}
 |\mathcal G_{g,h}(0;0)|
 \leq C_h(2\pi)^{-2g}N!R^{-N}N^{-5/2}.
\end{equation}
The bound is absolute; it does not require positivity of the
transformed coefficients.
\end{prp}
\begin{prf}
Choose a small fixed $r_0>0$ for which
Proposition~\ref{prp:b:common-domains} applies on $[-r_0,0]$.
On this interval put $S=R-x$. After scaling $y=S Y$, its first
two derivative bounds give the uniform hypotheses of
Lemma~\ref{lmm:b:coefficient-bound}. Thus
\[
 |L_{g,h}(0;x)|\leq C_h(2\pi)^{-2g}N!
                       (R-x)^{-N}N^{-3/2}.
\]
Since $N\geq3$,
$\int_{-r_0}^0(R-x)^{-N}\,\dd x\leq C R^{-N}/N$.
The near-zero part of \eqref{eq:b:primitive} has the bound
\eqref{eq:b:primitive-bound}.

\smallskip\noindent
\emph{The bounded negative interval.}
Let $u_{r_0}>0$ satisfy $w(-u_{r_0})=-r_0$ and choose
$U>u_{r_0}$. Lemma~\ref{lmm:b:negative-disks} gives a uniform
inverse radius strictly larger than $R$ for
$u\in[u_{r_0},U]$. Its lower bound increases with $u$, so
this radius may be chosen independently of $U$. Proposition~\ref{prp:b:tube} makes the
actual Borel functions holomorphic and bounded on a common circle
of radius $R_1>R$. Cauchy's estimate bounds this integral, after
division by $(2\pi)^{-2g}N!$, by $C_hR_1^{-N}$.

\smallskip\noindent
\emph{Negative infinity.}
Increase $U$ until $c_h\sqrt u I_0(2\sqrt u)>2R_1$ for every
$u\geq U$. By Lemma~\ref{lmm:b:growing-disks}, the remaining
integral after the same division is bounded by
\[
 C_h\int_U^\infty I_0(2\sqrt u)^h
              [c_h\sqrt u I_0(2\sqrt u)]^{-N}\,\dd u.
\]
Because $N\geq h+2$, reserve $h+2$ powers of the bracket to make
the integral finite. The reserved integrand is a constant times
$u^{-(h+2)/2}I_0(2\sqrt u)^{-2}$, which is integrable by
\eqref{eq:b:modified-bounds}. Bound every remaining power by
$(2R_1)^{-1}$. This proves an $O_h(R_1^{-N})$ bound for this
part as well. Its exponential advantage over $R^{-N}$ absorbs
any fixed power of $N$. Combining the three pieces with
Lemma~\ref{lmm:b:primitive} proves the proposition, including all
finitely many small genera by enlarging $C_h$.
\end{prf}
